\section{Pulse Code Modulation (PCM)}
Il segnale analogico viene campionato a una frequenza di campionamento che deve essere almeno il doppio della frequenza massima del segnale (teorema di Nyquist), e quindi quantizzato in un certo numero di livelli, e quindi codificato in un certo numero di bit per ogni campione.\\
Il PCM è un metodo di codifica digitale che viene utilizzato per rappresentare i segnali analogici in forma digitale, e quindi permette di trasmettere i segnali su reti digitali, e quindi migliorare la qualità del segnale e aumentare la velocità di trasmissione, e supporta anche la trasmissione bidirezionale, e quindi permette di fornire connettività wireless ad alta velocità in ambienti in cui le comunicazioni in fibra ottica non sono pratiche.\\

Per il segnale voce la banda è a 4 kHz, quindi per il teorema di Nyquist la frequenza di campionamento deve essere almeno 8 kHz, e quindi con una quantizzazione a 8 bit per campione si ottiene una velocità di trasmissione di 64 kbps, si trasmette un campione ogni 125 µs (1/8000).\\

Quantizzazione - Ogni campione viene rappresentato da un numero binario di 8 bit, e quindi si ottiene una rappresentazione digitale del segnale analogico, e quindi permette di trasmettere i segnali su reti digitali, e quindi migliorare la qualità del segnale e aumentare la velocità di trasmissione, e supporta anche la trasmissione bidirezionale, e quindi permette di fornire connettività wireless ad alta velocità in ambienti in cui le comunicazioni in fibra ottica non sono pratiche.\\

Il segnale digitale viene interllacciato con i segnali provenienti da altri utenti per formare una Digital Hierarchy. Questa tecnica si chiama TDM (Time Division Multiplexing).\\

\section{Plesiochronous Digital Hierarchy (PDH)}
Il PDH è una tecnica di multiplexing che consente di combinare più segnali digitali a velocità diverse in un unico flusso di dati, e quindi permette di trasmettere i segnali su reti digitali, e quindi migliorare la qualità del segnale e aumentare la velocità di trasmissione, e supporta anche la trasmissione bidirezionale, e quindi permette di fornire connettività wireless ad alta velocità in ambienti in cui le comunicazioni in fibra ottica non sono pratiche.\\

Plesiocrane significa che il clock non era lo stesso. 

\section{Synchronous Digital Hierarchy (SDH) - Synchronous Optical Network (SONET)}
Fu standardizzata intorno agli '90, in italia dalla telecom. Siamo passati dalle reti PDH a quelle SDH/SONET per aggiungere controllo delle rete. LA rete è costituita da una serie di trasmettitori che li mette insieme dal multiplexing per creare il segnale di trasporto. Lungo la rete abbiamo degli amplificatori. ADM è un nodo.
STS DEMUX separa poi il segnale.\\

Questi 3 livelli di overhead mi permette un controllo del flusso, mi permette di aumentare la velocità in fibra. TUtti gli elementi della rete usano lo stesso clock, sono sincronizzati in ampiezza e in fase. Avere questa struttura ben organizzata ci permette di aprire il mio switch a determinati istanti di tempo ed estrarre il flusso di dati che mi interessa.\\

Ho informazioni sullo stato della rete, posso fare il monitoraggio, posso fare il controllo del flusso, posso fare la gestione degli errori.\\

Introduco quello che si chiama la qualità del servizio. Quindi l'operatore oltre a garantirti la banda, ti garantisce anche la qualità di servizio.\\

Abbiamo una struttura stratificata. Prendiamo il nostro segnale elettrico, viene inserito il mio dato e aggiungo un overhead di cammino, dove questo overhead viene letto solo ricezione. Quando faccio il multiplexing ci metto l'overhead di linea, quindi mi controlla la sincronizzazione e l'errore nella comunicazione. Poi aggiungo overhead di sezione, e infine trasformo il segnale elettrico in ottico (photonic).\\

Tre tipi di overhead:
\begin{itemize}
    \item Di cammino
    \item Di linea
    \item Di sessione
\end{itemize}

\paragraph{SONET: Frame STS-n}
